%*************************************
% ภาคผนวก ค.
%*************************************
\vspace{-0.5in}
% \vspace*{-0.5in}

\appendixChapterTitle{ก}{การติดตั้งและใช้งาน Cloud-Based Platform}
\section{การติดตั้งและใช้งาน Cloud-Based Platform}
\hspace*{1.3em}
ในการติดตั้งและใช้งาน Cloud-Based Platform สำหรับการสนับสนุนกิจกรรมการเรียนการสอนและกิจกรรมทางวิชาการ ผู้ใช้งานจำเป็นต้องเตรียมความพร้อมของระบบและทำการติดตั้งตามขั้นตอนดังนี้

\vspace{0.5em}
\textbf{ข้อกำหนดของระบบ (System Requirements)}
\vspace{0.5em}

\begin{mycustomenum2}
    \item \textbf{ซอฟต์แวร์ที่จำเป็น:}
    \begin{mycustomenum2}
        \item Bun 1.2.0 หรือสูงกว่า (\href{https://bun.sh/}{https://bun.sh/})
        \item Docker และ Docker Compose (\href{https://www.docker.com/}{https://www.docker.com/})
        \item PostgreSQL 15 หรือสูงกว่า
        \item RabbitMQ 3.12 หรือสูงกว่า
        \item Proxmox VE (สำหรับการใช้งานจริง)
    \end{mycustomenum2}

    \item \textbf{ข้อกำหนดของฮาร์ดแวร์:}
    \begin{mycustomenum2}
        \item RAM: อย่างน้อย 8 GB (แนะนำ 16 GB หรือมากกว่า)
        \item CPU: อย่างน้อย 4 cores
        \item พื้นที่เก็บข้อมูล: อย่างน้อย 50 GB
        \item เชื่อมต่ออินเทอร์เน็ตเสถียร
    \end{mycustomenum2}
\end{mycustomenum2}

\vspace{0.5em}
\textbf{การดาวน์โหลดและติดตั้งโปรเจ็กต์}
\vspace{0.5em}

\begin{mycustomenum2}
    \item เปิด Terminal หรือ Command Prompt และทำการโคลนโปรเจ็กต์จาก GitHub Repository โดยใช้คำสั่งดังภาพที่ \ref{figC:GitClone}

\begin{figure}[htbp]
\centering
\adjustbox{frame, width=0.9\textwidth}{
    \begin{minipage}{0.85\textwidth}
        \texttt{\footnotesize
        git clone https://github.com/akikungz/cloud-based-platform.git\\
        cd cloud-based-platform
        }
    \end{minipage}
}
\caption{\fontSixTeen{คำสั่งสำหรับโคลนโปรเจ็กต์}}
\label{figC:GitClone}
\end{figure}

    \item ติดตั้ง Dependencies ของโปรเจ็กต์ โดยใช้คำสั่งดังภาพที่ \ref{figC:BunInstall}

\begin{figure}[htbp]
\centering
\adjustbox{frame, width=0.9\textwidth}{
    \begin{minipage}{0.85\textwidth}
        \texttt{\footnotesize
        bun install
        }
    \end{minipage}
}
\caption{\fontSixTeen{คำสั่งสำหรับติดตั้ง Dependencies}}
\label{figC:BunInstall}
\end{figure}

    \item ตั้งค่า Environment Variables โดยคัดลอกไฟล์ตัวอย่าง โดยใช้คำสั่งดังภาพที่ \ref{figC:EnvSetup}

\begin{figure}[htbp]
\centering
\adjustbox{frame, width=0.9\textwidth}{
    \begin{minipage}{0.85\textwidth}
        \texttt{\footnotesize
        cp .env.example .env
        }
    \end{minipage}
}
\caption{\fontSixTeen{การตั้งค่า Environment Variables}}
\label{figC:EnvSetup}
\end{figure}
\end{mycustomenum2}

\section{การเริ่มต้นใช้งานระบบ}
\hspace*{1.3em}
หลังจากติดตั้งโปรเจ็กต์เรียบร้อยแล้ว สามารถเริ่มต้นใช้งานระบบได้ดังนี้

\vspace{0.5em}
\textbf{การรันระบบในโหมด Development}
\vspace{0.5em}

\begin{mycustomenum2}
    \item เริ่มต้น Infrastructure Services ด้วย Docker Compose โดยใช้คำสั่งดังภาพที่ \ref{figC:DockerUp}

\begin{figure}[htbp]
\centering
\adjustbox{frame, width=0.9\textwidth}{
    \begin{minipage}{0.85\textwidth}
        \texttt{\footnotesize
        docker-compose up -d
        }
    \end{minipage}
}
\caption{\fontSixTeen{การเริ่มต้น Infrastructure Services}}
\label{figC:DockerUp}
\end{figure}

    \item เริ่มต้น Midori (Frontend Application) โดยใช้คำสั่งดังภาพที่ \ref{figC:MidoriDev}

\begin{figure}[htbp]
\centering
\adjustbox{frame, width=0.9\textwidth}{
    \begin{minipage}{0.85\textwidth}
        \texttt{\footnotesize
        cd apps/midori\\
        bun dev
        }
    \end{minipage}
}
\caption{\fontSixTeen{การเริ่มต้น Midori Frontend}}
\label{figC:MidoriDev}
\end{figure}

    \item เริ่มต้น Momoi (API Service) ใน Terminal ใหม่ โดยใช้คำสั่งดังภาพที่ \ref{figC:MomoiDev}

\begin{figure}[htbp]
\centering
\adjustbox{frame, width=0.9\textwidth}{
    \begin{minipage}{0.85\textwidth}
        \texttt{\footnotesize
        cd apps/momoi\\
        bun dev
        }
    \end{minipage}
}
\caption{\fontSixTeen{การเริ่มต้น Momoi API Service}}
\label{figC:MomoiDev}
\end{figure}

    \item เริ่มต้น Yuzu (VM Operations Service) ใน Terminal ใหม่ โดยใช้คำสั่งดังภาพที่ \ref{figC:YuzuDev}

\begin{figure}[htbp]
\centering
\adjustbox{frame, width=0.9\textwidth}{
    \begin{minipage}{0.85\textwidth}
        \texttt{\footnotesize
        cd apps/yuzu\\
        bun dev
        }
    \end{minipage}
}
\caption{\fontSixTeen{การเริ่มต้น Yuzu VM Operations Service}}
\label{figC:YuzuDev}
\end{figure}

    \item เข้าใช้งานระบบผ่าน Web Browser ที่ \url{http://localhost:3000}
\end{mycustomenum2}

\section{การใช้งานส่วนประกอบต่าง ๆ ของระบบ}
\hspace*{1.3em}
ระบบ Cloud-Based Platform ประกอบด้วยส่วนประกอบหลัก 3 ส่วน ดังนี้

\vspace{0.5em}
\textbf{Midori - Frontend Application}
\vspace{0.5em}

\begin{mycustomenum2}
    \item เป็น Web Application ที่พัฒนาด้วย Next.js 15.5
    \item ใช้ Material UI 7.3 และ Tailwind CSS 4 สำหรับส่วน User Interface
    \item รองรับการใช้งานบนอุปกรณ์ Desktop และ Mobile
    \item มีฟีเจอร์หลัก ดังนี้:
    \begin{mycustomenum2}
        \item การจัดการ Virtual Machines (สร้าง, ติดตาม, จัดการ)
        \item Dashboard สำหรับดูการใช้งานทรัพยากร
        \item ระบบ Authentication และการจัดการบัญชีผู้ใช้
        \item รองรับการใช้งานแบบ Responsive Design
    \end{mycustomenum2}
    \item URL สำหรับการเข้าใช้งาน: \url{http://localhost:3000}
\end{mycustomenum2}

\vspace{0.5em}
\textbf{Momoi - API Service}
\vspace{0.5em}

\begin{mycustomenum2}
    \item เป็น Core API Service ที่พัฒนาด้วย Elysia.js
    \item ทำหน้าที่เป็นตัวกลางระหว่าง Frontend และ VM Operations Service
    \item มีฟีเจอร์หลัก ดังนี้:
    \begin{mycustomenum2}
        \item RESTful API สำหรับการทำงานของแพลตฟอร์ม
        \item OpenAPI Documentation แบบอัตโนมัติ
        \item การติดตาม Telemetry ด้วย Jaeger และ Prometheus
        \item การเชื่อมต่อกับ RabbitMQ สำหรับการประมวลผลแบบ Asynchronous
    \end{mycustomenum2}
    \item URL สำหรับ API: \url{http://localhost:3001}
\end{mycustomenum2}

\vspace{0.5em}
\textbf{Yuzu - VM Operations Service}
\vspace{0.5em}

\begin{mycustomenum2}
    \item เป็น Service สำหรับจัดการ Virtual Machines ใน Proxmox VE
    \item ประมวลผลคำขอผ่าน RabbitMQ Message Queue
    \item มีฟีเจอร์หลัก ดังนี้:
    \begin{mycustomenum2}
        \item จัดการ VM Lifecycle (สร้าง, เริ่ม, หยุด, ลบ VMs)
        \item การ Provisioning แบบ Template-based
        \item การจัดการ IP Address แบบอัตโนมัติ
        \item ระบบ Logging ที่ครอบคลุมด้วย Pino
    \end{mycustomenum2}
\end{mycustomenum2}

\section{การตรวจสอบและติดตามระบบ}
\hspace*{1.3em}
ระบบมี Services สำหรับการตรวจสอบและติดตามการทำงานดังนี้

\vspace{0.5em}
\textbf{Monitoring Services}
\vspace{0.5em}

\begin{mycustomenum2}
    \item \textbf{Prometheus} - การเก็บข้อมูล Metrics
    \begin{mycustomenum2}
        \item URL: \url{http://localhost:9090}
        \item ใช้สำหรับติดตามประสิทธิภาพของระบบ
    \end{mycustomenum2}

    \item \textbf{Jaeger} - การติดตาม Distributed Tracing
    \begin{mycustomenum2}
        \item URL: \url{http://localhost:16686}
        \item ใช้สำหรับติดตามการทำงานของ API calls
    \end{mycustomenum2}

    \item \textbf{Grafana} - Dashboard สำหรับการแสดงผล
    \begin{mycustomenum2}
        \item URL: \url{http://localhost:3002}
        \item Username/Password: admin/admin (ค่าเริ่มต้น)
        \item ใช้สำหรับแสดงข้อมูลสถิติและ Metrics
    \end{mycustomenum2}

    \item \textbf{RabbitMQ Management} - การจัดการ Message Queue
    \begin{mycustomenum2}
        \item URL: \url{http://localhost:15672}
        \item Username/Password: user/password (ค่าเริ่มต้น)
        \item ใช้สำหรับตรวจสอบ Message Queue
    \end{mycustomenum2}
\end{mycustomenum2}

\section{การแก้ไขปัญหาเบื้องต้น}
\hspace*{1.3em}
หากพบปัญหาในการใช้งาน สามารถแก้ไขได้ดังนี้

\vspace{0.5em}
\textbf{ปัญหาที่พบบ่อยและวิธีแก้ไข}
\vspace{0.5em}

\begin{mycustomenum2}
    \item \textbf{Port ถูกใช้งานแล้ว}
    \begin{mycustomenum2}
        \item ตรวจสอบ Port ที่ใช้งานด้วยคำสั่ง: lsof -i :3000 (สำหรับ macOS/Linux)
        \item หยุดกระบวนการที่ใช้ Port หรือเปลี่ยน Port ในไฟล์ .env
    \end{mycustomenum2}

    \item \textbf{Docker Services ไม่สามารถเริ่มได้}
    \begin{mycustomenum2}
        \item ตรวจสอบสถานะด้วยคำสั่ง: docker-compose ps
        \item ดู Log ด้วยคำสั่ง: docker-compose logs [service\_name]
        \item รีสตาร์ท Services ด้วยคำสั่ง: docker-compose restart
    \end{mycustomenum2}

    \item \textbf{Database Connection Error}
    \begin{mycustomenum2}
        \item ตรวจสอบว่า PostgreSQL Container ทำงานอยู่
        \item ตรวจสอบการตั้งค่าใน Environment Variables
        \item รันคำสั่ง Database Migration (หากจำเป็น)
    \end{mycustomenum2}

    \item \textbf{Frontend ไม่สามารถเชื่อมต่อ API ได้}
    \begin{mycustomenum2}
        \item ตรวจสอบว่า Momoi API Service ทำงานอยู่ที่ Port 3000
        \item ตรวจสอบการตั้งค่า API URL ในไฟล์ Environment
        \item ตรวจสอบ CORS Settings ใน API Service
    \end{mycustomenum2}
\end{mycustomenum2}

% \newpage
\vspace{0.5em}
\textbf{การเข้าถึงข้อมูลเพิ่มเติม}
\vspace{0.5em}

\begin{mycustomenum2}
    \item อ่าน Documentation เพิ่มเติมใน README.md ของแต่ละ Component
    \item ตรวจสอบ Log Files สำหรับข้อมูลการ Debug
    \item ใช้ Jaeger UI สำหรับติดตามการทำงานของ API
    \item ใช้ Grafana Dashboard สำหรับติดตามประสิทธิภาพระบบ
\end{mycustomenum2}
